% Generated draft: requires mathematical and source review.
\documentclass[11pt]{article}
\usepackage[margin=25mm]{geometry}
\usepackage{fontspec}
\setmainfont{Latin Modern Roman}
\usepackage{amsmath,amssymb,mathtools,unicode-math}
\setmathfont{Latin Modern Math}
\usepackage{enumitem,needspace,xurl,hyperref}
\setlength{\parindent}{0pt}
\setlength{\parskip}{5pt}
\setlength{\emergencystretch}{3em}
\newcommand{\N}{\mathbb{N}}
\newcommand{\Z}{\mathbb{Z}}
\newcommand{\Q}{\mathbb{Q}}
\newcommand{\R}{\mathbb{R}}
\newcommand{\C}{\mathbb{C}}
\newcommand{\F}{\mathbb{F}}
\newcommand{\A}{\mathbb{A}}
\newcommand{\PP}{\mathbb{P}}
\newcommand{\E}{\mathbb{E}}
\newcommand{\eps}{\varepsilon}
\newcommand{\dd}{\,\mathrm{d}}
\begin{document}
\section*{UnsolvedMath ID: 30004525}
\textbf{Source problem code:} catalogue entry
\subsection{Definitions and mathematical statement}
In the usual commutative convention, an algebra \(A\) is Artinian if its ideals satisfy the descending chain condition. It is Gorenstein if, for every maximal ideal \(\mathfrak m\), the socle
\[
\operatorname{Soc}(A_{\mathfrak m})=\{z\in A_{\mathfrak m}:\mathfrak m A_{\mathfrak m}z=0\}
\]
is one-dimensional over the residue field \(A_{\mathfrak m}/\mathfrak m A_{\mathfrak m}\). An ideal is principal if it has the form \((u)=Au\). The source does not specify a base field, its characteristic, whether the algebra is local, or the grading assumptions for all parts of the question.

\textbf{Requested constructions.} In the intended class of Artinian Gorenstein algebras, the supplied record asks the following questions.
\begin{enumerate}
\item Given principal ideals \(J=(u)\) and \(K=(v)\), can the double \(Q_{J,K}\)-stratification be redefined so that its pieces have the expected dimensions and it retains its relation to the single-ideal stratifications \(Q_J(a)\) and \(Q_K(a)\)?
\item For nongraded Artinian Gorenstein algebras, can the \(Q_{\ell,m}\)-decomposition be used to define a Jordan order type?
\item For a binomial, can the \(Q_{x,y}\)-decomposition be determined by combinatorial invariants of its associated poset? Can such a description also cover variables of different degrees and the case of three variables?
\end{enumerate}

\textbf{Unresolved formulation.} The supplied record does not define any of these \(Q\)-stratifications or decompositions, the index \(a\), the expected dimensions, or the relations that a redefinition must preserve. It also does not define Jordan order type, specify the roles or admissible choices of \(\ell,m,x,y\), or construct an algebra and an associated poset from a binomial. Thus these are source-attributed construction questions; the supplied information does not determine a precise existence assertion or a combinatorial determination theorem. In particular, no choice of grading, binomial construction, poset, or dimension formula is imposed here.
\subsection{Short English statement}
The source asks for compatible double stratifications for two principal ideals, a Jordan order type for nongraded Artinian Gorenstein algebras, and a combinatorial description of the binomial case, including unequal variable degrees and three variables. The defining constructions and target conditions are missing, so a fully precise mathematical statement requires further source information.
\subsection{Sources}
\begin{itemize}
\item Local UnsolvedMath registry: permanent ID \(30004525\), title \emph{Double Jordan Stratifications of Artinian Gorenstein Algebras}, source problem code \texttt{catalogue entry}, in the Oberwolfach Reports Open Problems collection. The supplied proposed year is 2020; no individual proposer is identified.
\item Supplied report citation: \emph{Lefschetz Properties in Algebra, Geometry and Combinatorics} (2021), Oberwolfach Reports, pp.\ 1539--1600, DOI \texttt{10.4171/owr/2020/31}. \url{https://doi.org/10.4171/owr/2020/31}. The background describes Oberwolfach Reports as documenting mathematical workshops and their open-problem discussions. The literature-source list labels this reference \emph{Oberwolfach Reports 2020/31}.
\item Supplied literature reference: Iarrobino--Macias Marques--McDaniel, \emph{Jordan Type}. \url{https://arxiv.org/abs/1811.02072}.
\item The supplied literature assessment, dated exactly \texttt{2026-08-22}, labels the questions open and states that the proposed generalizations and combinatorial descriptions remain open, while existing work establishes only special Jordan-string decompositions. This is supplied dated curation derived from the OWR record, not an independently verified literature assessment; its status and cited claims require independent verification.
\item The registry supplies the statement-status label \texttt{corrected\_verified} and says that mathematical content was preserved while notation, typography, omitted definitions, or duplicate clauses were repaired for self-contained publication. This is a supplied editorial assertion; the record reproduced here still omits definitions needed for a precise formulation.
\end{itemize}
\end{document}
